\section{方法论：Agentic Engineering vs. Vibe Coding}

Peter 在访谈中给出了一组高传播度的对照：\textbf{Agentic Engineering} 与 \textbf{Vibe Coding}。这不是词汇游戏，而是工程纪律问题。

\subsection{为什么“vibe coding 是贬义”}

他将 vibe coding 描述为“凌晨三点后的模式”：高速度、低约束、第二天要花大量时间清理。相对地，agentic engineering 强调前置约束与可验证变更。

\begin{importantbox}{Agentic Engineering 的最小纪律}
\begin{itemize}
    \item 先定义目标和边界，再放权给 agent；
    \item 优先让 agent 读代码和文档，而不是先改；
    \item 用触发词控制模式切换（讨论模式 / 执行模式）；
    \item 每轮执行后检查副作用，而非只看功能是否“跑通”。
\end{itemize}
\end{importantbox}

\subsection{语音驱动开发：效率与人体工学的交换}

访谈里他提到大量使用 voice input，甚至出现“短期失声”。这提示我们：交互摩擦降低了，但新的工作负担（持续口述、上下文管理）会上升。

\begin{warningbox}{Voice-first 不等于总是更高效}
语音交互在 ideation 阶段有效，但在高精度约束任务中，文本依然更易审计和追踪。建议将 voice 用于“意图表达”，将关键约束落回文本。
\end{warningbox}

\subsection{本章小结}
Agentic Engineering 的核心不是“更会用模型”，而是\textbf{让模型在可治理边界内加速工程}。速度本身不是目标，可持续迭代才是目标。

